Een kaduuk zandkasteel.
Een gekunsteld toneel.

Een houten tol rolde.
Een lappenpop tolde

rondom een kinderbeen.
Speelgoed op het hoogveen.

Waardevolle schatten.
Een kind kon bevatten

wat spelen kon schenken
zonder na te denken

over gewenst profijt.
Het was de peutertijd,

waarin Emmy nog was
zonder ego poespas.

Ze deelde haar spullen
met spruitjes en knullen

zonder enige zorg.
Gift, lening of een borg

kwam nog niet in haar op.
Het was toch maar een pop.

Ze wilde geen baten,
maar het liefst speelmaten.

Vader had het gespaard.
Zijn werktijden verzwaard.

Dus moeder greep hard in.
Ze uitte haar weerzin

dat ze samen speelden
en haar poppen deelde.

“Emmy,” zei haar moeder,
als gefaald opvoeder,

terwijl ze ruw rukte
uit een diep bedrukte

buurmeisje haar handen,
de losse haarbanden

van een gammele pop.
Emmy keek naar de kop

en mama zag haar fout.
Emmy werd nu echt stout

zag ze aan het gezicht
van het volwassen wicht.

"Die pop was heel erg duur.
Het is dan heel erg zuur

als je die pop uitleent.”
Haar vriendin keek versteend,

toen ook de rest verdween.
Moeder zei toen gemeen;

“Dat haar ouders arm zijn,
dat is niet onze pijn.

Laat ze het zelf klaren
door ook te gaan sparen.
Heb je al je speelgoed?
Kijk nog eens extra goed!”

Emmy keek ma strak aan
en wilde niet meegaan.

Het buurmeisje keek rond,
greep de tol van de grond

en gaf het bang aan ma.
Het was een dilemma

voor de jonge Emmy.
Ze wilde affectie

van zowel haar moeder
als elke speelbroeder.

Hoewel nog zo piepjong,
doorzag ze de oorsprong

van haar moeders zorgen.
De geldangst voor morgen

als wel de achterdocht,
die een uitlaatklep zocht.

Zo werd moeder nukkig
en dolongellukig.

Toch werd er een zaadje
geplant in een laatje

in het nog jonge brein.
Ze moest oplettend zijn.

Ze raakte zo verhard.
Ze raakte ook verward.

Er was niets kwijtgeraakt
en/of kapotgemaakt.

En mocht dat wel spelen,
kon haar dat niets schelen.

Emma was haar gezwel.
Zij zag het als een spel

om Emmy’s weldaden
almaar te verraden.

“Mama, die lappenpop,
met dat oog uit zijn kop,

lag bij het buurmeisje.”
Status was haar prijsje.

Boekhouden haar hobby.
De balans van Emmy

was altijd negatief.
Emma gaf narratief

de verhoudingen weer.
Zij deed altijd veel meer

voor het dorp of gezin
dan ooit de benjamin

tot dusver presteerde,
was wat ze beweerde.

Zij kreeg alle ethiek,
Emmy alle kritiek

wat ze vaak negeerde.
Zus Emma lanceerde

de bombardementen
juist op de momenten

dat Emmy fideel was
en goedig in haar sas

in haar eentje speelde.
De vals veroordeelde

kon zich niet wapenen
en haar ontwapenen.

Emmy was afwezig.
Zij hield zich niet bezig

met moraalboekhouden.
Emma tegenhouden

lukte haar daarom niet.
Emmy nam haar verdriet

of werd ontzettend kwaad.
De agressieve daad

van de boze Emmy
gaf Emma munitie.
